\chapter{Events and Probability}

\paragraph{Exercise 1.1} We flip a fair coin ten times.
\begin{enumerate}
  \item[(a)] Our sample space $\Omega$ is the set of all possible outcomes of ten
    fair coin flips. All simple events have equal probability. The sample space has
    $2^{10}$ simple events, so each of them has probability $2^{-10}$. We are
    interested in the probability of the event $E_a$ that the number of heads and
    the number of tails are equal; that is, in the event that there are 5 heads.
    There are $\binom{10}{5}$ simple events that fulfil this condition. Consequently,
    \[
      \pr(E_a) = \binom{10}{5} \cdot 2^{-10} = 252 \cdot 2^{-10} = \frac{63}{256}.
    \]

  \item[(b)] We are interested in the probability of the event $E_b$ that there are
    more heads than tails. Let $E_i$ be the event that the number of heads equals
    $i \in [0, 10]$. Then $E = \bigcup_{i = 6}^{10} E_i$. Since $E_1, E_2,..., E_{10}$
    are mutually disjoint,
    \begin{align*}
      \pr\left(E_b\right)
        &= \pr\left(\bigcup_{i = 6}^{10} E_i \right)      = \sum_{i=6}^{10} \pr\left(E_i\right) = 2^{-10} \sum_{i=6}^{10} \binom{10}{i} \\
        &= 2^{-10} \left( 210 + 120 + 45 + 10 + 1\right)  = \frac{193}{512}.
    \end{align*}

  \item[(c)] We are interested in the probability of the event $E_c$ that the $i$th
    flip and the $(11 - i)$th flip are the same for $i \in [1,5]$. Let
    $X_i \in$ \{H, T\} be the outcome of the $i$th flip. Since we are tossing a
    fair coin, $\pr\left(X_i = \text{H} \right) = \pr\left(X_i = \text{T} \right) = \frac{1}{2}$.
    For $X_{11-i}$ holds, $\pr\left(X_{11-i} = X_i \right) = \pr\left(X_{11-i}
    \not= X_i \right) = \frac{1}{2}$. Since the outcome of each toss is independet
    of the outcome of the other nine tosses,
    \[
      \pr\left(E_c\right)
        = \pr\left(\bigcap_{i=1}^5 X_i = X_{11-i}\right)
        = 2^{-5}.
    \]

  \item[(d)] We are interested in the probablity of the event $E_d$ that we flip
    at least four consecutive heads. Let $i \in [1,10]$ and let $E_i$ be the event
    that we flip $i$ consecutive heads. Then,
    \[ \pr\left(E_d\right) = \pr\left(\bigcup_{i=4}^{10} E_i\right).\]
    We have to consider, that the events $E_1,..., E_{10}$ are not mutually disjoint.
    For instance, the sequence (H,H,H,H,T,H,H,H,H,H) is element both of the event
    $E_4$ and $E_5$. Let's have a look at the individual events $E_4,..., E_{10}$
    there $x \in$ \{H,T\}:
    \begin{align*}
      \begin{tabular}{lll}
        $E_4$ &=  &\{(H,H,H,H,T,x,x,x,x,x), (T,H,H,H,H,T,x,x,x,x), (x,T,H,H,H,H,T,x,x,x), \\
              &   &\ \,(x,x,T,H,H,H,H,T,x,x), (x,x,x,T,H,H,H,H,T,x), (x,x,x,x,T,H,H,H,H,T), \\
              &   &\ \,(x,x,x,x,x,T,H,H,H,H)\} \\
        $E_5$ &= &\{(H,H,H,H,H,T,x,x,x,x), (T,H,H,H,H,H,T,x,x,x), (x,T,H,H,H,H,H,T,x,x), \\
              &   &\ \,(x,x,T,H,H,H,H,H,T,x), (x,x,x,T,H,H,H,H,H,T), (x,x,x,x,T,H,H,H,H,H)\} \\
        $E_6$ &=  &\{(H,H,H,H,H,H,T,x,x,x), (T,H,H,H,H,H,H,T,x,x), (x,T,H,H,H,H,H,H,T,x), \\
              &   &\ \,(x,x,T,H,H,H,H,H,H,T), (x,x,x,T,H,H,H,H,H,H)\} \\
        $E_7$ &=  &\{(H,H,H,H,H,H,H,T,x,x), (T,H,H,H,H,H,H,H,T,x), (x,T,H,H,H,H,H,H,H,T), \\
              &   &\ \,(x,x,T,H,H,H,H,H,H,H)\} \\
        $E_8$ &=  &\{(H,H,H,H,H,H,H,H,T,x), (T,H,H,H,H,H,H,H,H,T), (x,T,H,H,H,H,H,H,H,H)\} \\
        $E_9$ &=  &\{(H,H,H,H,H,H,H,H,H,T), (T,H,H,H,H,H,H,H,H,H)\} \\
        $E_{10}$ &=  &\{(H,H,H,H,H,H,H,H,H,H)\}.
      \end{tabular}
    \end{align*}
    Now we can determine the cardinality of $E_4, ..., E_{10}$:
    \begin{align*}
      \begin{tabular}{ll}
        $|E_4| = 2 \cdot 2^5 + 5 \cdot 2^4 = 144$   &$|E_8| = 2 \cdot 2^1 + 1 \cdot 2^0 = 5$ \\
        $|E_5| = 2 \cdot 2^4 + 4 \cdot 2^3 = 64$    &$|E_9| = 2 \cdot 2^0 = 2$ \\
        $|E_6| = 2 \cdot 2^3 + 3 \cdot 2^2 = 28$    &$|E_{10}| = 1.$ \\
        $|E_7| = 2 \cdot 2^2 + 2 \cdot 2^1 = 12$    & \\
      \end{tabular}
    \end{align*}
    The events $E_4 \cup E_5, E_6, E_7, E_8, E_9, E_{10}$ are mutually disjoint.
    Therefore,
    \[ \pr\left(E_d\right) = \pr\left(\bigcup_{i=4}^{10} E_i\right)
      = \pr\left(E_4 \cup E_5 \right) + \sum_{i=6}^{10} \pr\left(E_i\right).\]
    It holds that
    \[
      \pr\left(E_4 \cup E_5 \right)
      = \pr\left(E_4\right) + \pr\left(E_5 \right) + \pr\left(E_4 \cap E_5 \right)
      = 2^{-10} \cdot 144 + 2^{-10} \cdot 64 - 2^{-10} \cdot 2
      = \frac{103}{512}.
    \]
    Now, we can compute the probability for event $E_d$:
    \[ \pr\left(E_d\right)
      = \pr\left(E_4 \cup E_5 \right) + \sum_{i=6}^{10} \pr\left(E_i\right) \\
      = \frac{103}{512} + 2^{-10}(28 + 12 + 5 + 1) \\
      = \frac{63}{256}.
    \]
\end{enumerate}

\paragraph{Exercise 1.2} Let $i \in \{ 1,2 \}$ and $X_i$ be the outcome of roll $i$.
\begin{enumerate}
  \item[(a)] We are interested in the probability of the event that $X_1 = X_2$,
    that is
    \[
      \pr(X_1 = X_2) = \sum_{i=1}^6 \pr(X_1 = i \cap X_2 = i).
    \]
    Since the outcomes of the rolls are independent,
    \begin{align*}
      \pr(X_1 = X_2)  &= \sum_{i=1}^6 \pr(X_1 = i \cap X_2 = i) \\
                      &= \sum_{i=1}^6 \pr(X_1 = i) \cdot \pr(X_2 = i) \\
                      &= 6 \cdot \frac{1}{6} \cdot \frac{1}{6} \\
                      &= \frac{1}{6}.
    \end{align*}

  \item[(b)] We are interested in the probability of the event that $X_1 > X_2$,
    that is
    \begin{align*}
      \pr(X_1 > X_2)  &= \sum_{i=1}^6 \pr(X_1 = i \cap X_2 < i) \\
                      &= \sum_{i=1}^6 \pr(X_1 = i) \cdot \pr(X_2 < i) \\
                      &= \frac{1}{6} \sum_{i=1}^6 \pr(X_2 < i) \\
                      &= \frac{1}{6} \sum_{i=1}^6 (i-1) \cdot \frac{1}{6} \\
                      &= \frac{1}{36} \sum_{i=1}^5 i.
    \end{align*}
    Now we can apply the Small Gauss sum formula to get
    \[
      \pr(X_1 > X_2) = \frac{1}{36} \sum_{i=1}^5 i
        = \frac{1}{36} \cdot \frac{5 \cdot (5 + 1)}{2}
        = \frac{5}{12}.
    \]

  \item[(c)] We are interested in the probability of the event that $X_1 + X_2$
    is even. The sum of $X_1$ and $X_2$ is even if and only if either both $X_1$
    and $X_2$ are even or both $X_1$ and $X_2$ are odd,
    \[
      \pr(X_1 + X_2 = 0 \text{  mod }2) = \pr\left(X_1 = X_2 \text{  mod }2\right).
    \]
    Since the outcomes of the rolls are independent and  the events $X_i = 0 \text{  mod }2$ and $X_i = 1 \text{  mod }2)$ are disjoint, for $i \in \{ 1,2 \}$,
    \begin{align*}
      \pr(X_1 + X_2 = 0 \text{  mod }2)
         &= \pr((X_1 = 0 \text{  mod }2 \cap X_2 = 0 \text{  mod }2)  \quad \cup \quad (X_1 = 1 \text{  mod }2 \cap X_2 = 1 \text{  mod }2)) \\
         &= \pr(X_1 = 0 \text{  mod }2 \cap X_2 = 0 \text{  mod }2)     \quad  +   \quad \pr(X_1 = 1 \text{  mod }2 \cap X_2 = 1 \text{  mod }2) \\
         &= \pr(X_1 = 0 \text{  mod }2) \cdot \pr(X_2 = 0 \text{  mod }2)\quad  +  \quad \pr(X_1 = 1 \text{  mod }2) \cdot \pr(X_2 = 1 \text{  mod }2) \\
         &= \frac{1}{2} \cdot \frac{1}{2} \quad  +   \quad \frac{1}{2} \cdot \frac{1}{2} \\
         &= \frac{1}{2}.
  \end{align*}

  \item[(d)] The value of the product of the dice is in the interval [1,36]. The
    set $S$ of all perfect squares in this interval equals $\{ 1, 4, 9, 16, 25, 36 \}$.
    So we are interested in the probability that the product of the dice is an
    element of $S$, $\pr\left(X_1 \cdot X_2 \in S\right)$. Then one has
    \[
      \pr\left(X_1 \cdot X_2 \in S\right)
        = \sum_{s \in S}\pr\left(X_1 \cdot X_2 = s\right)
        = \frac{1}{36} + \frac{3}{36} +  \frac{1}{36} +  \frac{1}{36} +  \frac{1}{36} +  \frac{1}{36}\\
        = \frac{2}{9}.
    \]
\end{enumerate}

\paragraph{Exercise 1.3} We shuffle a standard deck of cards, obtaining a permutation
that is uniform over all 52! possible permutations.
\begin{enumerate}
  \item[(a)] We are interested in the probability of the event $E_a$ that the
    first two cards include at least one ace. Therefore we compute the complementary
    probability of $E_a$, that is the probability $\pr\left(\bar{E_a}\right)$
    that the first two cards inlude no ace. Since a standard deck of cards
    includes 4 aces,
    \[ \pr\left(\bar{E_a}\right) = \frac{48}{52} \cdot \frac{47}{51} = \frac{188}{221}.\]
    Consequently,
    \[ \pr\left(E_a\right) = 1 - \pr\left(\bar{E_a}\right) = 1 - \frac{188}{221} = \frac{33}{221}.\]

  \item[(b)] Let $E_b$ be the event that the first five cards include at least one
    ace. Again we compute the complementary probability of the specified event.
    \[
      \pr\left(E_b\right)
        = 1 - \pr\left(\bar{E_b}\right)
        = 1 - \frac{48}{52} \cdot \frac{47}{51} \cdot \frac{46}{50} \cdot \frac{45}{49} \cdot \frac{44}{48}
        = 1 - \frac{205,476,480}{311,875,200}.
    \]
    Alternatively, we could argue that the described situation suites the conditions
    for a hypergeometric situation, so that
    \[
      \pr\left(E_b\right)
        = 1 - \frac{\binom{4}{0} \cdot \binom{48}{5}}{\binom{52}{5}}
        = 1 - \frac{\binom{48}{5}}{\binom{52}{5}}.
    \]

  \item[(c)] Let $E_c$ be the event that the first two cards are a pair of the
    same rank.  $E_c$ occurs if the second card has the same rank as the first
    card. Therefore, one has
    \[ \pr\left(E_c\right) = \frac{3}{51} = \frac{1}{17}.\]

  \item[(d)] Let $E_d$ be the event that the first five cards are all diamonds. Then,
    \[
      \pr\left(E_d\right)
        = \frac{\binom{13}{5} \cdot \binom{39}{0}}{\binom{52}{5}}
        = \frac{\binom{13}{5}}{\binom{52}{5}}
        = \frac{13 \cdot 12 \cdot 11 \cdot 10 \cdot 9}{52 \cdot 51 \cdot 50 \cdot 49 \cdot 18}
        = \frac{33}{66,640}.
    \]

  \item[(e)] Let $E_e$ be the event that the first five cards form a full house
    (three of one rank and two of another rank). Let us think of event $E_e$ as
    drawing 5 cards from a deck. Two of them have to be from different rank, so
    there are $\binom{13}{2}$ possibilities two choose two different ranks.
    Altogether we have to draw three from one of the picked rank and two of the
    other. Since it doesn't matter which rank appears three and which two times,
    one has
    \[
      \pr\left(E_e\right)
        = \frac{\binom{13}{2} \cdot \binom{4}{3} \cdot \binom{4}{2} \cdot 2}{\binom{52}{5}}
        = \frac{78 \cdot 4 \cdot 6 \cdot 2}{2,598,960}
        = \frac{6}{4,165}.
    \]
\end{enumerate}

\paragraph{Exercise 1.4} Let us play a tournament in which we stop as soon as one
of us wins n games. We are evenly matched, so each of us wins any game with
probability $\frac{1}{2}$, independently of other games. What is the probability
that the loser has won $k$ games when the match is over? \\[0.2cm]
Let $X$ be the number of games the loser has won. We are interested in the
probability that $X = k$, that is the probability that the tournament ends after
$n + k$ games. Therefore, the winner has to win $n-1$ games among the first
$n + k - 1$ games and additionally has to win the $(n+k)$-th game. Every tournament
that fulfills this condition has probability $\frac{1}{2}^n \cdot \frac{1}{2}^{k}$.
Since there are $\binom{n+k-1}{n-1}$ possibilities to choose $n-1$ games the winner
wins among $n+k-1$ games,
\[
  \pr\left(X = k\right)
    = \binom{n+k-1}{n-1} \cdot \left(\frac{1}{2}\right)^n \cdot \left(\frac{1}{2}\right)^{k}
    = \binom{n+k-1}{n-1} \cdot \frac{1}{2^{n+k}}.
\]
To verify this result, we show that the following equation holds.
\[
  \sum_{k=0}^{n-1}\pr\left(X=k\right) = 1.
\]
Because of the symmetry of the binomial coefficient, that is $\binom{b}{a} = \binom{b}{b - a}$,
\begin{align*}
  \sum_{k=0}^{n-1}\pr\left(X=k\right)
    &= \sum_{k=0}^{n-1} \binom{n+k-1}{n-1} \cdot \frac{1}{2^n} \cdot \frac{1}{2^k} \\
    &= \sum_{k=0}^{n-1} \binom{n+k-1}{k} \cdot \frac{1}{2^n} \cdot \frac{1}{2^k} \\
    &= \frac{1}{2^n} \cdot \sum_{k=0}^{n-1} \binom{n+k-1}{k} \cdot \frac{1}{2^k}.
\end{align*}
Since $\binom{b}{a} = (-1)^a \binom{a - b - 1}{a}$,
\begin{align*}
  \sum_{k=0}^{n-1}\pr\left(X=k\right)
    &= \frac{1}{2^n} \cdot \sum_{k=0}^{n-1} \binom{n+k-1}{k} \cdot \frac{1}{2^k} \\
    &= \frac{1}{2^n} \cdot \sum_{k=0}^{n-1} \binom{k+n-1}{k} \cdot (-1)^k \left(- \frac{1}{2}\right)^{k} \\
    &= \frac{1}{2^n} \cdot \sum_{k=0}^{n-1} \binom{-n}{k} \cdot \left(- \frac{1}{2}\right)^{k}.
\end{align*}
Under consideration of the binomial series $(1 + x)^b = \sum_{a=0}^{\infty}
\binom{b}{a}x^a, \quad |x| < 1, b \in \mathbb{R}$, we finally get
\begin{align*}
  \sum_{k=0}^{n-1}\pr\left(X=k\right)
    &= \frac{1}{2^n} \cdot \sum_{k=0}^{n-1} \binom{-n}{k} \cdot \left(- \frac{1}{2}\right)^{k} \\
    &= \frac{1}{2^n} \cdot \left(1 - \frac{1}{2}\right)^{-n} \\
    &= \frac{1}{2^n} \cdot 2^n \\
    &= 1.
\end{align*}

\paragraph{Exercise 1.5}
\begin{enumerate}
  \item[(a)] Assume that Bob chooses die A and Alice chooses die B. In the
    following we list all of the possible events and their probabilities.
    \begin{align*}
    \begin{tabular}{c|c|c|c}
        die A   &die B  &Probability    &Winner \\
        \hline
            1       &2      &1/3 $\cdot$ 1/3 = 1/9 &Alice\\
            1       &4      &1/9        &Alice\\
            1       &9      &1/9        &Alice\\
            6       &2      &1/9        &Bob\\
            6       &4      &1/9        &Bob\\
            6       &9      &1/9        &Alice\\
            8       &2      &1/9        &Bob\\
            8       &4      &1/9        &Bob\\
            8       &9      &1/9        &Alice\\
    \end{tabular}
  \end{align*}
    Therefore the probability that Alice wins is $5 \cdot \frac{1}{9} =
    \frac{5}{9} > \frac{1}{2}$.

  \item[(b)] Assume that Bob chooses die B and Alice chooses die C. In the
    following we list all of the possible events and their probabilities.
    \begin{align*}
    \begin{tabular}{c|c|c|c}
        die B   &die C  &Probability  &Winner \\
        \hline
            2       &3      &1/9      &Alice\\
            2       &5      &1/9      &Alice\\
            2       &7      &1/9      &Alice\\
            4       &3      &1/9      &Bob\\
            4       &5      &1/9      &Alice\\
            4       &7      &1/9      &Alice\\
            9       &3      &1/9      &Bob\\
            9       &5      &1/9      &Bob\\
            9       &7      &1/9      &Bob\\
    \end{tabular}
  \end{align*}
    Again, the probability that Alice wins is $5 \cdot \frac{1}{9} =
    \frac{5}{9} > \frac{1}{2}$.

  \item[(c)] Assume that Bob chooses die C and Alice chooses die A. In the
    following we list all of the possible events and their probabilities.
    \begin{align*}
    \begin{tabular}{c|c|c|c}
        die C   &die A  &Probability  &Winner \\
        \hline
            3       &1      &1/9      &Bob\\
            3       &6      &1/9      &Alice\\
            3       &8      &1/9      &Alice\\
            5       &1      &1/9      &Bob\\
            5       &6      &1/9      &Alice\\
            5       &8      &1/9      &Alice\\
            7       &1      &1/9      &Bob\\
            7       &6      &1/9      &Bob\\
            7       &8      &1/9      &Alice\\
    \end{tabular}
  \end{align*}
    And still, the probability that Alice wins is $5 \cdot \frac{1}{9} =
    \frac{5}{9} > \frac{1}{2}$.

\end{enumerate}

\paragraph{Exercise 1.6} Balls-and-bin game: We start with one black ball and one
white ball in a bin. We repeatedly choose one ball from the bin uniformly at
random, and then put the ball back in the bin with another ball of the same color.
We repeat until there are $n$ balls in the bin. \\[0,2cm]
Let $X$ be the number of white balls and $k \in \mathbb{N}$, so that $0 \leq k
\leq n-2$. We want to show that,
\[ \pr(X=k+1) = \frac{1}{n-1}. \]
To model this $(n-2)$-stage balls-and-bin game, we describe the drawing of a white
ball as the event 1 and and a black ball as the event 0 respectively. Then our
sample space $\Omega$ can be defined as $\Omega := \Omega_1 \times ... \times
\Omega_{n-2}$, where $\Omega_j = \{ 0, 1\}$ and $j = 1,...,n-2$. Now one concrete
balls-and-bin game can be described as an elementary event $\omega =
(a_1,...,a_{n-2}) \in \Omega$, where $a_j \in \Omega_j$ specifies the $j$-th ball
that we draw. Then
\begin{align*}
  \pr(\omega)
    &= \pr(a_{n-2}| a_1,...,a_{n-3}) \cdot \pr(a_1,...,a_{n-3}) \\
    &= \pr(a_1) \cdot \pr(a_2| a_1) \cdot \pr(a_3| a_1,a_2) \cdot ... \cdot \pr(a_{n-2}| a_1,...,a_{n-3}).
\end{align*}
Let $p_j(1)$ be the probability that the $j$-th ball that we draw is
a white ball and $p_j(0)$ the probability that the $j$-th ball is a black one. \\
Initially there is one white and one black ball in the bin. Therefore,
\[ p_1(1) = \frac{1}{2} = p_1(0). \]
Suppose we have drawn $j-1$ balls and there has been $w$ white and $j-1-w$ black
ones among them. Consequently there are $1 + w$ white balls and $1 + j - 1 - w$
black balls in the bin before the draw the $j$-th ball. Let
$\left(a_1,...,a_{j-1}\right)$ be a tuple with $w$ 1's and $(j-1-w)$ 0's, that is
$\sum_{i=1}^{j-1}a_i = w$. Then
\begin{align*}
    p_j(1|a_1,...,a_{j-1}) = \frac{1 + w}{1 + 1 + (j-1)}, \\
    p_j(0|a_1,...,a_{j-1}) = \frac{1+j-1-w}{1 + 1 + (j-1)}.
\end{align*}
Let $\omega = (a_1, ..., a_{n-2}) \in \Omega$, so that $\sum_{j=1}^{n-2}a_j = k$
and $0 \leq k \leq n-2$. Then
\[
  \pr(\omega) = \frac {\prod_{j=0}^{k-1}(1 + j) \prod_{j=0}^{(n-2)-k-1}(1 + j)}
                      {\prod_{j=0}^{n-3}(1+1+j)}.
\]
Since the event $X=k+1$ consists of all tuples $(a_1,...,a_{n-2})$ with
$\sum_{j=1}^{n-2}a_j = k$ and since each of these tuples has the above probability,
\begin{align*}
  \pr(X=k+1)
    &= \binom{n-2}{k} \cdot \frac {\prod_{j=0}^{k-1}(1 + j) \prod_{j=0}^{(n-2)-k-1}(1 + j)}{\prod_{j=0}^{n-3}(1+1+j)} \\
    &= \frac{(n-2)!}{k! (n-2-k)!} \cdot \frac{k! \cdot (n-2-k)!}{(n-1)!} \\
    &= \frac{1}{n-1}.
\end{align*}
Hence the number of white balls is equally likely to be any number between 1 and
$n-1$.

\paragraph{Exercise 1.7}
\begin{enumerate}
  \item[(a)] \textbf{Lemma 1.3} (\textcolor{WildStrawberry}{Inclusion-exclusion
    principle}): Let $E_1,...,E_n$ be any $n$ events. Then
    \[
      \pr\left(\bigcup_{i=1}^n E_i\right)
      = \sum_{i=1}^n \pr(E_i) - \sum_{i<j}\pr(E_i \cap E_j) + \sum_{i<j<k}\pr(E_i \cap E_j \cap E_k) - ... + (-1)^{l+1} \sum_{i_1<i_2< ... <i_l} \pr\left(\bigcap_{r=1}^lE_{i_r}\right)+ ... .
    \]
    \textit{Proof}. We want to show the Inclusion-exclusion principle by induction
    over $n \in \mathbb{N}$, where $n \geq 2$. \\
    For the base case (IB) assume that $n = 2$. According to Lemma 1.1, one has
    \begin{align*}
      \pr\left(E_1 \cup E_2 \right)
        &= \pr(E_1) + \pr(E_2) - \pr\left(E_1 \cap E_2 \right) \\
        &= \sum_{i=1}^2 \pr(E_i) - \sum_{i<j}\pr\left(E_1 \cap E_2 \right).
    \end{align*}
    Let $n \geq 2$, so that the Inclusion-exclusion principle holds for $n$
    (induction hypothesis, IH). We have to show, that the Inclusion-exclusion
    principle then also hold for $n+1$.
    \begin{align*}
      \pr\left(\bigcup_{i=1}^{n+1} E_i\right)
        &= \pr\left(\left(\bigcup_{i=1}^{n} E_i\right) \cup E_{n+1}\right) \\
        &=_{\text{Lemma 1.1}} \pr\left(\bigcup_{i=1}^{n} E_i\right) + \pr(E_{n+1}) -
            \pr\left(\left(\bigcup_{i=1}^{n} E_i\right) \cap E_{n+1}\right) \\
        &= \pr\left(\bigcup_{i=1}^{n} E_i\right) + \pr(E_{n+1}) -
            \pr\left(\bigcup_{i=1}^{n} \left(E_i \cap E_{n+1}\right)\right) \\
        &=_{IH} \pr(E_{n+1}) +
          \left(
            \sum_{i=1}^n \pr(E_i) - \sum_{i<j}\pr(E_i \cap E_j) + ... + (-1)^{n+1}
            \sum_{i_1<i_2< ... <i_n} \pr\left(\bigcap_{r=1}^n E_{i_r}\right)
          \right) \\
        &\ \ - \Bigg(
            \sum_{i=1}^n \pr\left(E_i \cap E_{n+1}\right) -
            \sum_{i<j} \pr\left((E_i \cap E_{n+1}) \cap (E_j \cap E_{n+1})\right)
            + ... \\
        &\quad \quad \quad+ (-1)^{n+1} \sum_{i_1<i_2< ... <i_n} \pr\left(\bigcap_{r=1}^n
            \left(E_{i_r} \cap E_{n+1}\right)\right)
          \Bigg) \\
        &= \sum_{i=1}^{n+1} \pr(E_i) - \sum_{i<j}\pr(E_i \cap E_j) + ... +
          (-1)^{n+1} \sum_{i_1<i_2< ... <i_n} \pr\left(\bigcap_{r=1}^n E_{i_r}\right) \\
        &\ \ - \Bigg(
            \sum_{i=1}^n \pr\left(E_i \cap E_{n+1}\right) - \sum_{i<j}\pr\left(E_i
            \cap E_j \cap E_{n+1}\right) + ...
         + (-1)^{n+1} \sum_{i_1<i_2< ... <i_n} \pr\left(\bigcap_{r=1}^n
          \left(E_{i_r} \cap E_{n+1}\right)\right)
          \Bigg) \\
        &= \sum_{i=1}^{n+1} \pr(E_i) - \sum_{i<j}\pr(E_i \cap E_j) +  ... + (-1)^{n+2}
          \sum_{i_1<i_2< ... <i_{n+1}} \pr\left(\bigcap_{r=1}^{n+1}E_{i_r}\right).
    \end{align*}
    Hence, the Inclusion-exclusion principle applies.

  \item[(b)] \textbf{Lemma 1.8}: Let $E_1,...,E_n$ be any $n$ events and $l \in
  \mathbb{N}$, so that $l$ is odd. Then
  \[
    \pr\left(\bigcup_{i=1}^n E_i\right)
    \leq \sum_{i=1}^n \pr(E_i) - \sum_{i<j}\pr(E_i \cap E_j) + \sum_{i<j<k}\pr(E_i \cap E_j \cap E_k) - ... + (-1)^{l+1} \sum_{i_1<i_2< ... <i_l} \pr\left(\bigcap_{r=1}^lE_{i_r}\right).
  \]
  \textit{Proof}. Let $l \in \mathbb{N}$, so that $3 \leq l \leq n$ and $l$ is odd. According to Lemma 1.3 one has,
  \begin{align*}
  \begin{tabular}{l}
    $\pr\left(\bigcup_{i=1}^n E_i\right) -
    \left(
      (-1)^{l+2} \sum_{i_1<i_2< ... <i_{l+1}} \pr\left(\bigcap_{r=1}^{l+1} E_{i_r}\right) + ... +
      (-1)^{n+1} \pr\left(\bigcap_{i=1}^n E_{i}\right)
    \right)$ \\
    $\ \ = \sum_{i=1}^n \pr(E_i) + ... + (-1)^{l+1} \sum_{i_1<i_2< ... <i_l}
      \pr\left(\bigcap_{r=1}^lE_{i_r}\right).$
  \end{tabular}
  \end{align*}
  Therefore, we have to show that
  \[
    (-1)^{l+2} \sum_{i_1<i_2< ... <i_{l+1}} \pr\left(\bigcap_{r=1}^{l+1} E_{i_r}\right)
    + ... + (-1)^{n+1} \pr\left(\bigcap_{i=1}^n E_{i}\right) \leq 0.
  \]
  Since $l$ is odd,
  \begin{align*}
      &(-1)^{l+2} \sum_{i_1<i_2< ... <i_{l+1}} \pr\left(\bigcap_{r=1}^{l+1} E_{i_r}\right) + ... + (-1)^{n+1} \pr\left(\bigcap_{i=1}^n E_{i}\right) \\
    = &-\sum_{i_1<i_2< ... <i_{l+1}} \pr\left(\bigcap_{r=1}^{l+1} E_{i_r}\right)
      + ... + (-1)^{n+1} \pr\left(\bigcap_{i=1}^n E_{i}\right).
  \end{align*}
  Furthermore, for any $1 \leq k < n$ holds,
  \[
    \sum_{i_1<i_2< ... <i_k} \pr\left(\bigcap_{r=1}^k E_{i_r}\right)
    \geq \sum_{i_1<i_2< ... <i_{k+1}} \pr\left(\bigcap_{r=1}^{k+1} E_{i_r}\right).
  \]
  Hence, for any pair $(k,k+1)$, there $l+1 \leq k < n$ and $k$ is even, holds
  \begin{align*}
    \   &(-1)^{k+1} \sum_{i_1<i_2< ... <i_{k}} \pr\left(\bigcap_{r=1}^{l+1} E_{i_r}\right) + (-1)^{k+2} \pr\left( \bigcap_{r=1}^{k+1} E_{i} \right) \\
    =   &-\sum_{i_1<i_2< ... <i_k} \pr\left( \bigcap_{r=1}^{l+1} E_{i_r} \right) + \pr\left(\bigcap_{r=1}^{k+1} E_{i}\right) \\
    \leq &\ 0.
  \end{align*}
  Therefore,
  \[
    (-1)^{l+2} \sum_{i_1<i_2< ... <i_{l+1}} \pr\left(\bigcap_{r=1}^{l+1} E_{i_r}\right)
    + ... + (-1)^{n+1} \pr\left(\bigcap_{i=1}^n E_{i}\right) \leq 0,
  \]
  and
  \[
    \pr\left(\bigcup_{i=1}^n E_i\right)
    \leq \sum_{i=1}^n \pr(E_i) - \sum_{i<j}\pr(E_i \cap E_j) + \sum_{i<j<k}\pr(E_i \cap E_j \cap E_k) - ... + (-1)^{l+1} \sum_{i_1<i_2< ... <i_l} \pr\left(\bigcap_{r=1}^lE_{i_r}\right).
  \]

  \item[(c)] \textbf{Lemma 1.9}: Let $E_1,...,E_n$ be any $n$ events and $l \in
  \mathbb{N}$, so that $l$ is even. Then
  \[
    \pr\left(\bigcup_{i=1}^n E_i\right)
    \geq \sum_{i=1}^n \pr(E_i) - \sum_{i<j}\pr(E_i \cap E_j) + \sum_{i<j<k}\pr(E_i \cap E_j \cap E_k) - ... + (-1)^{l+1} \sum_{i_1<i_2< ... <i_l} \pr\left(\bigcap_{r=1}^lE_{i_r}\right).
  \]
  \textit{Proof}. Let $l \in \mathbb{N}$, so that $2 \leq l \leq n$ and $l$ is even. Similarly to exercise 1.7b) we want now to show that
  \[
    (-1)^{l+2} \sum_{i_1<i_2< ... <i_{l+1}} \pr\left(\bigcap_{r=1}^{l+1} E_{i_r}\right)
    + ... + (-1)^{n+1} \pr\left(\bigcap_{i=1}^n E_{i}\right) \geq 0.
  \]
  Since $l$ is even,
  \begin{align*}
      &(-1)^{l+2} \sum_{i_1<i_2< ... <i_{l+1}} \pr\left(\bigcap_{r=1}^{l+1} E_{i_r}\right) + ... + (-1)^{n+1} \pr\left(\bigcap_{i=1}^n E_{i}\right) \\
    = &\sum_{i_1<i_2< ... <i_{l+1}} \pr\left(\bigcap_{r=1}^{l+1} E_{i_r}\right)
      - ... + (-1)^{n+1} \pr\left(\bigcap_{i=1}^n E_{i}\right).
  \end{align*}
  For any pair $(k,k+1)$, there $l+1 \leq k < n$ and $k$ is odd, holds
  \begin{align*}
    \   &(-1)^{k+1} \sum_{i_1<i_2< ... <i_{k}} \pr\left(\bigcap_{r=1}^{l+1} E_{i_r}\right) + (-1)^{k+2} \pr\left( \bigcap_{r=1}^{k+1} E_{i} \right) \\
    =   &\sum_{i_1<i_2< ... <i_k} \pr\left( \bigcap_{r=1}^{l+1} E_{i_r} \right) - \pr\left(\bigcap_{r=1}^{k+1} E_{i}\right) \\
    \geq &\ 0.
  \end{align*}
  Therefore,
  \[
    (-1)^{l+2} \sum_{i_1<i_2< ... <i_{l+1}} \pr\left(\bigcap_{r=1}^{l+1} E_{i_r}\right)
    + ... + (-1)^{n+1} \pr\left(\bigcap_{i=1}^n E_{i}\right) \geq 0,
  \]
  and
  \[
    \pr\left(\bigcup_{i=1}^n E_i\right)
    \geq \sum_{i=1}^n \pr(E_i) - \sum_{i<j}\pr(E_i \cap E_j) + \sum_{i<j<k}\pr(E_i \cap E_j \cap E_k) - ... + (-1)^{l+1} \sum_{i_1<i_2< ... <i_l} \pr\left(\bigcap_{r=1}^lE_{i_r}\right).
  \]

\end{enumerate}

\paragraph{Exercise 1.8} Let us choose a number uniformly at random from the range
$[1, 1,000,000]$. Using the inclusion–exclusion principle, we want to determine
the probability that the number chosen is divisible by one or more of 4, 6, and 9. \\
Let $E_i$ be the event that the choosen number is divisible by $i\in \mathbb{N}^+$.
Starting from 1, every $i$-th natural number is divisible by $i$. Therefore
\begin{align*}
  \pr\left(E_4\right) &= \frac{\left\lfloor \frac{1,000,001}{4} \right\rfloor}{1,000,001} = \frac{250,000}{1,000,001} \\
  \pr\left(E_6\right) &= \frac{\left\lfloor \frac{1,000,001}{6} \right\rfloor}{1,000,001} = \frac{166,666}{1,000,001} \\
  \pr\left(E_9\right) &= \frac{\left\lfloor \frac{1,000,001}{9} \right\rfloor}{1,000,001} = \frac{111,111}{1,000,001}.
\end{align*}
A number $n \in \mathbb{N}$ is divisible by two numbers $i,j \in \mathbb{N}^+$
if ans only if $n$ is divisible by the least common multiple. The least
common multiple of 4 and 6 is 12, of 4 and 9 is 36, of 6 and 9 is 18 and of 4, 6
and 9 is 36. Consequently
\begin{align*}
  \pr\left(E_4 \cap E_6\right) &= \frac{\left\lfloor \frac{1,000,001}{12} \right\rfloor}{1,000,001} = \frac{83,333}{1,000,001} \\
  \pr\left(E_4 \cap E_9\right) &= \frac{\left\lfloor \frac{1,000,001}{36} \right\rfloor}{1,000,001} = \frac{27,777}{1,000,001} \\
  \pr\left(E_6 \cap E_9\right) &= \frac{\left\lfloor \frac{1,000,001}{18} \right\rfloor}{1,000,001} = \frac{55,555}{1,000,001} \\
  \pr\left(E_4 \cap E_6 \cap E_9\right) &= \frac{\left\lfloor \frac{1,000,001}{36} \right\rfloor}{1,000,001} = \frac{27,777}{1,000,001}.
\end{align*}
Applying the inclusion-exclusion principle, we are now able to determine the
probability that the chosen number is divisible by one or more of 4, 6, and 9.
\begin{align*}
  \pr\left(\bigcup_{i\in \{4,6,9\}}E_i\right)
    &= \sum_{i\in \{4,6,9\}} \pr\left(E_i\right) -
       \left(\pr\left(E_4 \cap E_6\right) + \pr\left(E_4 \cap E_9\right) + \pr\left(E_6 \cap E_9\right)\right) + \pr\left(E_4 \cap E_6 \cap E_9 \right) \\
    &= \frac{527,777}{1,000,001} - \frac{166,665}{1,000,001} + \frac{27,777}{1,000,001} \\
    &= \frac{388,889}{1,000,001}.
\end{align*}

\paragraph{Exercise 1.9} Suppose that a fair coin is flipped $n$ times. Let $i, k \in
\mathbb{N}^+, 1 \leq i \leq n - \left(\log_2(n) + k \right) + 1$ and let $E_i$ be the
event that there is a sequence of $\log_2(n) + k$ consecutive heads, starting with
the $i$th flip. \\
We are interested in an upper bound on $\pr\left(\bigcup_i E_i\right)$. Let's first
determine the probability $\pr(E_i)$ that a coin comes up heads $\log_2(n) + k$
times in succession. Since the coin is fair, one has
\[ \pr(E_i)
  = \left(\frac{1}{2}\right)^{\log_2(n) + k}
  = \frac{1}{2^{\log_2(n)}} \cdot \frac{1}{2^k}
  = \frac{1}{n} \cdot \frac{1}{2^k}
  = \frac{1}{2^k n}.
\]
Note that the events $E_1,...,E_{n - \left(\log_2(n) + k \right) + 1}$ are not
mutually disjoint. Therefore $\pr\left(\bigcup_i E_i\right) = \sum_i \pr(E_i)$
doesn't apply. However, since we are only interested in an upper bound on
$\pr\left(\bigcup_i E_i\right)$ we can apply the union bound (Lemma 1.2),
\[ \pr\left(\bigcup_i E_i\right) \leq \sum_i \pr(E_i). \]
Therefore
\[ \pr\left(\bigcup_i E_i\right)
  \leq \left(n - \left(\log_2(n) + k \right) + 1\right) \cdot \frac{1}{2^k n}
  = \frac{1}{2^k} - \frac{\log_2(n) + k - 1}{2^k n}
  \leq \frac{1}{2^k}.
\]
Hence $\frac{1}{2^k}$ is an upper bound on the probability that there is a
sequence of $\log_2(n) + k$ consecutive heads, with $k \in \mathbb{N}^+$.

\paragraph{Exercise 1.10} Assume that we have a fair coin and a two-headed coin.
We choose one of the two coins randomly with equal probability and flip it. Given
that the flip was heads, what is the probability that we flipped the two-headed coin? \\
Let $F$ be the event that we choose the fair coin and $U$ the event that we choose
the two-headed coin. Let $H$ be the event that the choosen coin comes up heads
and $T$ the event that it comes up tails. \\
We are interested in the conditional probability $\pr(U|H)$. Applying Bayes' Law
one has
\[ \pr(U|H)
    = \frac{\pr(U \cap H)}{\pr(H)}
    =\frac{\pr(H|U) \cdot \pr(U)}{\pr(H|U) \cdot \pr(U) + \pr(H|F) \cdot \pr(F)}
    =\frac{1 \cdot \frac{1}{2}}{1 \cdot \frac{1}{2} + \frac{1}{2} \cdot \frac{1}{2}}
    = \frac{2}{3}.
\]

\paragraph{Exercise 1.11} Assume we are trying to send a single bit, either a 0
 or a 1. When we transmit the bit, it goes through a series of $n$ relays
 before it arrives. Each relay flips the bit independently with probability $p$.
 \begin{enumerate}
   \item[(a)] We are interested in the probabilty to receive the correct bit.
    If an odd number of relays flips the bit an incorrect bit arrives and if
    an even number of relays flips the bit the correct bit arrives. Let $k \in
    \mathbb{N}$ and let $E_k$ be the event that exactly $2 \cdot k$ relays flip
    the bit. Then
    \[ \pr\left(\cup_{k=0}^{\left\lfloor n/2 \right\rfloor} E_k \right) \]
    is the probability to receive the correct bit.
    Since $E_0,...,E_{\left\lfloor n/2 \right\rfloor}$ are mutually
    disjoint events,
    \[
      \pr\left(\cup_{k=0}^{\left\lfloor n/2 \right\rfloor} E_k \right)
      = \sum_{k=0}^{\left\lfloor n/2 \right\rfloor} \pr\left( E_k \right).
    \]
    Let $k \in \mathbb{N}$. There are $\binom{n}{2k}$ possibilities to choose $2k$
    relays among the $n$ relays that flip the bit. The probability that $2k$
    specified relays flip the bit, is $p^{2k}(1-p)^{n-2k}$. Therefore
    \[ \pr\left(E_k\right) = \binom{n}{2k} p^{2k}(1-p)^{n-2k}, \]
    and
    \[
      \pr\left(\cup_{k=0}^{\left\lfloor n/2 \right\rfloor} E_k \right)
      = \sum_{k=0}^{\left\lfloor n/2 \right\rfloor} \pr\left( E_k \right)
      = \sum_{k=0}^{\left\lfloor n/2 \right\rfloor} \binom{n}{2k} p^{2k}(1-p)^{n-2k}.
    \]

   \item[(b)] Let us say the relay has bias $q$, there $q \in \mathbb{R}$,
   $q \in [-1, 1]$ if the probability it flips the bit is $(1 - q)/2$. We want
   to show that sending a bit through two relays with bias $q_1$ and $q_2$ is
   equivalent to sending a bit through a single relay with bias $q_1q_2$. \\
   Let $X \in $ \{correct, incorrect\} indicate whether we receive the correct bit
   after sending it through two relays with bias $q_1$ and $q_2$. Similarly, let
   $Y \in $ \{correct, incorrect\} indicate whether we receive the correct bit
   after sending it through one relays with bias $q_1q_2$. According to the
   definition of the bias of a relay, one has
   \[ \pr(Y = \text{incorrect}) = \frac{1 - q_1q_2}{2}. \]
   Let's consider the event $\{ X = \text{incorrect} \}$. This event occurs if and
   only if exactly one of the two relays flips the bit. Therefore
   \begin{align*}
     \pr\left( X = \text{incorrect} \right)
      &= \frac{1-q_1}{2} \cdot \left(1 - \frac{1-q_2}{2}\right) +
         \frac{1-q_2}{2} \cdot \left(1 - \frac{1-q_1}{2}\right) \\
      &= \frac{1-q_1}{2} - \frac{(1-q_1)(1-q_2)}{4} +
         \frac{1-q_2}{2} - \frac{(1-q_2)(1-q_1)}{4} \\
      &= \frac{1}{4} \left( 2(1-q_1) - (1-q_1)(1-q_2) + 2(1-q_2) - (1-q_2)(1-q_1) \right) \\
      &= \frac{1}{2} \left((1-q_1) - (1-q_1)(1-q_2) + (1-q_2) \right) \\
      &= \frac{1}{2} \left(1 - q_1 - (1 - q_2 - q_1 + q_1q_2) + 1 - q_2 \right) \\
      &= \frac{1}{2} \left(1 - q_1 - 1 + q_2 + q_1 - q_1q_2 + 1 - q_2 \right) \\
      &= \frac{1}{2} \left(1 - q_1q_2 \right). \\
   \end{align*}
   Hence, $\pr\left( X = \text{incorrect} \right) = \pr\left( Y = \text{incorrect}
   \right)$, and consequently sending a bit through two relays with bias $q_1$ and
   $q_2$ is equivalent to sending a bit through a single relay with bias $q_1q_2$.

   \item[(c)] Theorem: The probability you receive the correct bit when it passes
    through $n$ relays is
    \[ \frac{1 + (1 - 2p)^n}{2}. \]
    \textit{Proof}. We show the above theorem by induction over $n \in \mathbb{N}^+$.
    As the base case, assume that $n = 1$. The probability you receive the correct
    bit equals the probability that the only relay doesn't flip, that is $(1 - p)$.
    Since
    \[ \frac{1 + (1 - 2p)^1}{2} = \frac{2 - 2p}{2} = (1 - p), \]
    the theorem applies for the base case. \\
    Let $n \in \mathbb{N}^+$, so that the probability you receive the correct bit
    when it passes through $n$ relays is $\frac{1 + (1 - 2p)^n}{2}$ (induction
    hypothesis, IH). We have to show that the theorem also applies for $n+1$
    relays. Let $E_1$ be the event that the bit is correct after it passes through
    $n$ relays and the $(n+1)$th relay doesn't flip the bit. Let $E_2$ be the
    event that the bit is incorrect after it passes through $n$ relays and the
    $(n+1)$th relay flips the bit. Let $E$ be the event that you receive the correct
    bit when it passes through $n + 1$ relays. Then $E = E_1 \cup E_2$. Since
    $E_1$ and $E_2$ are mutually disjoint, one has
    \begin{align*}
      \pr(E)
      &= \pr(E_1) + \pr(E_2) \\
      &=_{\text{IH}}  \left( \frac{1 + (1 - 2p)^n}{2} \cdot (1 - p) \right) +
        \left( \left( 1 - \frac{1 + (1 - 2p)^n}{2} \right) \cdot p \right) \\
      &= \frac{1}{2} \left( (1-p) + (1-p)(1 - 2p)^n \right) +
        \frac{1}{2} \left( 2p - (p + p(1 - 2p)^n \right) \\
      &= \frac{1}{2} \left( 1 - p + (1-p)(1 - 2p)^n + 2p - p - p(1 - 2p)^n \right) \\
      &= \frac{1}{2} \left( 1  + (1 - 2p)^n(1 - p - p) \right) \\
      &= \frac{1  + (1 - 2p)^{n+1}}{2}.
    \end{align*}
    Hence, the probability you receive the correct bit when it passes through $n$
    relays is $\frac{1 + (1 - 2p)^n}{2}$.
 \end{enumerate}

\paragraph{Exercise 1.12} Monty Hall problem: \\
Intuitive solution: Assume that the contestant made her first choice. Since the
contestant had initially no further information , the probability that it is
behind one of the other curtains is $\frac{2}{3}$. Monty then opens one of the
other curtains to show a goat. Consequently the probability that the car is behind
the revealed curtain is 0. Since the probability that the car is behind one
of the non-chosen curtains is $\frac{2}{3}$, the probability that the car is behind
the not-chosen and not-revealed curtain is now $\frac{2}{3}$. Therefore, the
contestant should switch curtains.

Let's analyse the given problem formally. First, we define the sample space $\Omega$.
Let $i \in \{ 1, 2, 3 \}$ and $C_i$ the event that the car is behind curtain $i$.
Then $\Omega = \{ C_1, C_2, C_3 \}$. Furthermore, let $O_i$ be the event that Monty
opens curtain $i$. Without loss of generality assume that the contestant chooses
curtain 1 and that Monty opens curtain 2. Now we are interested in the probabilities
$\pr\left(C_1 | O_2\right)$ and $\pr\left(C_2 | O_2\right)$. Since $C_1, C_2$ and
$C_3$ are mutually disjoint sets such that $\cup_{i=1}^3 C_i = \Omega$, Bayes' Law
applies and one has
\[
  \pr\left(C_1 | O_2\right)
    = \frac {\pr\left(O_2 | C_1\right) \cdot \pr\left(C_1 \right)} {\sum_{i=1}^3 \pr\left(O_2 | C_i\right) \cdot \pr\left(C_i\right)}
    = \frac {\frac{1}{2} \cdot \frac{1}{3}}{\frac{1}{2} \cdot \frac{1}{3} + 0 \cdot \frac{1}{3} + 1 \cdot \frac{1}{3}}
    = \frac{1}{3},
\]
and
\[
  \pr\left(C_3 | O_2\right)
    = \frac {\pr\left(O_2 | C_3\right) \cdot \pr\left(C_3 \right)}{\sum_{i=1}^3 \pr\left(O_2 | C_i\right) \cdot \pr\left(C_i\right)}
    = \frac {1 \cdot \frac{1}{3}}{\frac{1}{2} \cdot \frac{1}{3} + 0 \cdot \frac{1}{3} + 1 \cdot \frac{1}{3}}
    = \frac{2}{3}.
\]
Since $\pr\left(C_3 | O_2\right) > \pr\left(C_1 | O_2\right)$, the contestant
should switch the curtains.

\paragraph{Exercise 1.13}  Let $D$ be the event that a person has the disorder and
$\bar{D}$ the event that the person has not the disorder. Our samples space $\Omega$
comprises the whole population. Obviously $D$ and $\bar{D}$ are mutually disjoint
sets such that $D \cup \bar{D} = \Omega$. Furthermore, let $P$ be the event that
a person is tested positive. We can now apply Bayes' Law to compute the probability
that a person chosen uniformely at random and tested positive has the disorder.
\[
  \pr(D|P)
    = \frac{\pr(P|D) \cdot \pr(D)}{\pr(P|D) \cdot \pr(D) + \pr(P|\bar{D}) \cdot \pr(\bar{D})}
    = \frac{0.999 \cdot 0.02}{0.999 \cdot 0.02 + 0.005 \cdot 0.98}
    = \frac{0.01998}{0.02488}
    = \frac{999}{1244}
    \approx 0.803.
\]
Hence, the probability that a positive tested person has the disorder is
approximately 0.803.

\paragraph{Exercise 1.14} I am playing in a racquetball tournament, and I am up
against a player I have watched but never played before. I consider three
possibilities for my prior model: we are equally talented, and each of us is equally
likely to win each game; I am slightly better, and therefore I win each game
independently with probability 0.6; or he is slightly better, and thus he wins
each game independently with probability 0.6. Before we play, I think that each
of these three possibilities is equally likely. In our match we play until one
player wins three games. I win the second game, but he wins the first, third, and
fourth. After this match, in my posterior model, with what probability should I
believe that my opponent is slightly better than I am?
Let $B$ be the event that I am slightly better than my opponent, $E$ the event that
we are equally talented, and $W$ the event that my opponent is slightly better than
I am. $B, E$ and $W$ are mutually disjoint events. Let LWLL be the event that I
win the second game and my opponent wins th first, third and fourth game. \\
We will apply Bayes' Law to compute the probability that my opponent is slightly
better than I am under the condition that LWLL occurs.
\begin{align*}
  \pr(W|\text{LWLL})
    &= \frac{\pr(\text{LWLL}|W) \cdot \pr(W)}{\pr(\text{LWLL}|B) \cdot \pr(B) + \pr(\text{LWLL}|E) \cdot \pr(E) + \pr(\text{LWLL}|W) \cdot \pr(W)} \\
    &= \frac{\left(\frac{3}{5}\right)^3 \cdot \frac{2}{5} \cdot \frac{1}{3}} {\left( \frac{2}{5} \right)^3 \cdot \frac{3}{5} \cdot \frac{1}{3} + \left( \frac{1}{2} \right) ^4 \cdot \frac{1}{3} + \left(\frac{3}{5}\right)^3 \cdot \frac{2}{5} \cdot \frac{1}{3}} \\
    &= \frac{864}{1873} \\
    &\approx 0.461.
\end{align*}
Hence, in my posterior model, the probability that my opponent is slightly better
than I am is approximately 0.461.

\paragraph{Exercise 1.15} Suppose that we roll ten standard six-sided dice.
Let $n \in \mathbb{N}^+$ and let $E_n$ be the event that the sum of $n$ rolled
standard six-sided dice is divisible by 6. \\
We claim that $\pr(E_n) = \frac{1}{6}$. We will prove that statement by induction
on $n \in \mathbb{N}^+$. As the base case assume that $n=1$. There is exactly one
possible outcome that is divisible by 6 - namely 6. Therefore $\pr(E_1) = \frac{1}{6}$.
Let $n \in \mathbb{N}^+$, so that $\pr(E_n) = \frac{1}{6}$ (induction hypothesis,
IH). Let's consider $\pr(E_{n+1})$. $E_{n+1}$ occurs either when $E_{n+1}$ occurs
and our $(n+1)$th roll shows a 6, or when $E_{n+1}$ doesn't occur but the sum of
the first $n$ rolled dice is equal to 6 minus the outcome of the $(n+1)$th roll
modulo 6. Whatever case occurs, there is exactly one possible outcome of the
$(n+1)$th rolled dice, so that the total sum is divisble by 6. Therefore
\[
  \pr(E_{n+1})
    = \frac{1}{6} \cdot \pr(E_n) + \frac{1}{6} \cdot \pr(\bar{E_n})
    =_{\text{IH}} \frac{1}{6} \cdot \frac{1}{6} + \frac{1}{6} \cdot \frac{5}{6}
    = \frac{1}{6}.
\]
Hence, $\pr(E_n) = \frac{1}{6}$ for all $n \in \mathbb{N}^+$. Consequently, the
probability that the sum of 10 rolled standard six-sided dice is divisible by
6 is $\frac{1}{6}$.

\paragraph{Exercise 1.16} Let's formalize the setting. Our sample space $\Omega$
consists of all possible sequences of rolls. Furthermore, let $E_1$ be the event
that all three dice show the same number on the first roll. Let $E_2$ be the event
that exactly two of the three dice show the same number on the first roll, and let
$E_3$ be the event that all three dice show different numbers on the first roll.
Obviously $E_1, E_2$ and $E_3$ are mutually disjoint events such that
$E_1 \cup E_2 \cup \E_3 = \Omega$. Additionally, let $W$ be the event that the
player wins the game.
\begin{enumerate}
  \item[(a)] The probability that all three dice show a specified number is
    $\frac{1}{6^3}$. Since there are six different numbers and it does not matter
    which of them the dice show as long as they all show the same number,
    \[ \pr(E_1) = 6 \cdot \frac{1}{6^3} = \frac{1}{36}.\]

  \item[(b)] The probability that two dice show a specified number and one die
    shows any other number is $\frac{1}{6^2} \cdot \frac{5}{6} = \frac{5}{6^3}$.
    Since there are six different numbers and it does not matter which two of the
    three dice show the same number,
    \[ \pr(E_2) = \binom{3}{2} \cdot 6 \cdot \frac{5}{6^3} = \frac{5}{12}.\]

  \item[(c)] Assume that exactly two of the three dice show the same number on
    the first roll. According to the player's strategy, the player re-rolls the
    die that does not match. The probability that she wins after the second roll
    is therefore $\frac{1}{6}$. Assume that she does not win after the second roll.
    Then again, she re-rolls the die that does not match and the probability that
    the die shows the requested number is $\frac{1}{6}$. Therefore, for the
    probability that the player wins, conditioned on exactly two of the three dice
    showing the same number on the first roll holds
    \[ \pr(W|E_2) =\frac{1}{6} + \frac{5}{6} \cdot \frac{1}{6} = \frac{11}{36}. \]

  \item[(d)] Let's first determine the probability that all dice show different
    numbers on the first roll. This happens when the second die shows a different
    number than the first one, and the third die shows a different number than the
    first and the sceond die. Therefore $\pr(E_3) = \frac{5}{6} \cdot \frac{4}{6}
    = \frac{5}{9}$. \\
    Furthermore, let us consider the probability that the player wins, conditioned
    on all three dice showing different numbers on the first roll. According to the
    player's strategy, the player re-rolls all dice. Since the rolls are independent,
    the probability that she wins after the second roll is $\frac{1}{36}$. The
    probability that all dice are different and she wins after the third roll is
    $\frac{5}{9} \cdot \frac{1}{36}$. The probability that exactly two dice show the
    same number and she wins after the third roll is $\frac{5}{12} \cdot \frac{1}{6}$.  Altogether, one has
    \[ \pr(W|E_3)
        = \frac{1}{36} + \frac{5}{9} \cdot \frac{1}{36} +
          \frac{5}{12} \cdot \frac{1}{6}
        = \frac{73}{648}.
    \]
    We can now apply the Law of total probability to compute the probability that
    the player wins the game.
    \[
      \pr(W)
        = \sum_{i=1}^3 \pr(W|E_i) \cdot \pr(E_i)
        = 1 \cdot \frac{1}{36} + \frac{11}{36} \cdot \frac{5}{12} + \frac{73}{648} \cdot \frac{5}{9}
        = \frac{2,539}{11,664}
        \approx 0.218.
    \]
\end{enumerate}

\paragraph{Exercise 1.17} Let's analyse the matrix multiplication algorithm of
section 1.3 under the assumption that we are working over the integers modulo $k$
for $k > 2$. \\
Our modified algorithm chooses a random vector $\bar{r} = (r_1, r_2, ..., r_n)
\in \{ 0, 1, ..., k-1 \} ^n$. It then computes $\textbf{AB}\bar{r}$ by first
computing $\textbf{B}\bar{r}$ and then $\textbf{A}(\textbf{B} \bar{r})$, and it
also computes $\textbf{C} \bar{r}$. If $\textbf{A}(\textbf{B}\bar{r}) \not=
\textbf{C}\bar{r}$, then $\textbf{AB} \not= \textbf{C}$. Otherwise, it returns
that $\textbf{AB} = \textbf{C}$. \\
The probability that the algorithm returns that $\textbf{AB} = \textbf{C}$ when
they are actually not equal is bounded by the following theorem. \\
\textbf{Theorem}: If $\textbf{AB} \not= \textbf{C}$ and if $\bar{r}$ is chosen
uniformly at random from $\{ 0, 1, ..., k-1 \} ^n$, then
\[ \pr(\textbf{AB}\bar{r}) = \textbf{C}\bar{r}) \leq \frac{1}{k}. \]
\textit{Proof}. Let $\textbf{D} = \textbf{AB} - \textbf{C} \not= 0$. Then
$\textbf{AB}\bar{r} = \textbf{C}\bar{r}$ implies that $\textbf{D}\bar{r} = 0$.
Since $\textbf{D} \not= 0$ it must have some nonzero entry; without loss of
generality, let that entry be $d_{11}$. \\
For $\textbf{D}\bar{r} = 0$, it must be the case that
\[ \sum_{j=1}^n d_{1j}r_j = 0 \]
or, equivalently,
\begin{align} r_1 = \frac{\sum_{j=2}^n d_{1j}r_j}{d_{11}}. \end{align}
Suppose that we choose the $r_j$ independently and uniformly at random from
$\{ 0, 1, ..., k-1 \}$ in order, from $r_n$ down to $r_1$. Choosing the $r_j$ in
this way is equivalent to choosing a vector $\bar{r}$ uniformly at random. Now
consider the situation just before $r_1$ is chosen. At this point, the right-hand
side of Eqn. (1.1) is determined, and there is at most one choice for $r_1$ that
will make the equality hold. Since there are $k$ choices for $r_1$, the equality
holds with probability at most $\frac{1}{k}$, and hence the probability that
$\textbf{AB}\bar{r} = \textbf{C}\bar{r}$ is at most $\frac{1}{k}$.

\paragraph{Exercise 1.18} Suppose we are given an integer $z \in \{ 0, ..., n-1 \}$
and want to output a value that equals $F(z)$ with probability at least
$\frac{1}{2}$. \\
We use a randomized algorithm that chooses a random number $x \in \{ 0, ..., n-1 \}$.
It then computes $y = (z - x) \text{ mod }n$ and afterwards uses the lookup table
to output the value $(F(x) + F(y)) \text{ mod }m$. \\
If the adversary has not changed the values of $x$ and $y$ we can be sure that
our algorithm outputs the correct value. The probability that the adversary has
changed neither $x$ nor $y$ is $\frac{4}{5} \cdot \frac{4}{5} = \frac{16}{25}$.
Since $\frac{16}{25} \geq \frac{1}{2}$, our algorithm outputs a value that equals
F(z) with probability at least $\frac{1}{2}$.

Suppose that we are allowed to repeat our initial algorithm three times. Let $i
\in \{ 1, 2, 3 \}$ and let $a_i$ be the output of the $i$th application of the
algorithm. We use the following strategy: If at least two outcomes match, we output
the corresponding value. Otherwise the output $a_1$.  \\
 Let $C$ be the event that we output the correct answer. To assess $\pr(C)$, we
consider the probability of the complementary event. The probability that we
output the wrong value is greater than zero if and only if there are at least
two wrong outputs among the three applications of our algorithm. The probability
that there are at least two wrong outputs is $\left(\frac{9}{25}\right)^2$.
Therefore
\[ \pr(C) = 1 - \pr(\bar{C})\geq 1 - \left(\frac{9}{25}\right)^2 = \frac{544}{625}.\]
Hence,the probability $\pr(C)$ that our enhanced algorithm returns the correct
answer is at least $0.8704$.

\paragraph{Exercise 1.19} We roll two standard six-sided dice. Let $i \in \{ 1, 2 \}$
and let $X_i \in \{ 1, 2, 3, 4, 5, 6 \}$ be the outcome of the $i$-th die. Let
$X = X_1 + X_2$. Then
\begin{align*}
  \pr(X = 2 | X_1 = 6) = 0 < \frac{1}{36} = \pr(X = 2); \\
  \pr(X_1 = 2 | X_2 = 2) = \frac{1}{6} = \pr(X_1 = 2); \\
  \pr(X_1 = 1 | X = 2) = 1 > \frac{1}{6} = \pr(X_1 = 1).
\end{align*}

\paragraph{Exercise 1.20} $~$ \\[0.2cm]
\textbf{Theorem 1.10}: If $E_1, E_2, ..., E_n$ are mutually independent, then so
are $\bar{E_1}, \bar{E_2}, ..., \bar{E_n}$. \\[0.2cm]
\textit{Proof}. Let $E_1, E_2, ..., E_n$ be mutually independent. Then, for any
subset $I \subseteq [1, n]$,
\begin{align*}
  \pr\left(\bigcap_{i \in I} \bar{E_i}\right)
    &= \pr\left(\bar{\bigcup_{i \in I} E_i}\right) \\
    &= 1 - \pr\left(\bigcup_{i \in I} E_i\right) \\
    &= 1 - \left( \sum_{i \in I}\pr(E_i) - \sum_{i,j \in I, i < j} \pr(E_i \cap E_j) + ... + (-1)^{|I| + 1}\pr\left(\bigcap_{i\in I}E_i\right) \right) \\
    &= 1 - \left( \sum_{i \in I} \pr(E_i) - \sum_{i,j \in I, i < j} \pr(E_i) \cdot \pr(E_j) + ... + (-1)^{|I| + 1}\prod_{i\in I} \pr(E_i) \right) \\
    &= 1 - \sum_{i \in I}\pr(E_i) + \sum_{i,j \in I, i < j} \pr(E_i) \cdot \pr(E_j) + ... + (-1)^{|I|}\prod_{i\in I} \pr(E_i) \\
    &= \prod_{i \in I}\left(1 - \pr(E_i)\right) \\
    &= \prod_{i \in I}\pr(\bar{E_i}).
\end{align*}
We have used De Morgan's law in the first line, the inclusion-exclusion principle
in the third line and the assumption that $E_1, E_2, ..., E_n$ are mutually
independent combined with definition 1.3 in the fourth line.

\paragraph{Exercise 1.21} We choose a number uniformly at random from the set
$\{ 1, 2, 3, 4 \}$. Then $\Omega = \{ 1, 2, 3, 4 \}$. Let $X, Y, Z$ be three events,
such that $X = \{ 1, 2 \}, Y = \{ 2, 3 \}$ and $Z = \{ 1, 3 \}$. Since
\begin{align*}
  \pr(X \cap Y) = \pr\left(\{ 2 \}\right) = \frac{1}{4}
    = \frac{1}{2} \cdot \frac{1}{2} = \pr(X) \cdot \pr(Y), \\
  \pr(X \cap Z) = \pr\left(\{ 1 \}\right) = \frac{1}{4}
    = \frac{1}{2} \cdot \frac{1}{2} = \pr(X) \cdot \pr(Z), \\
  \pr(Y \cap Z) = \pr\left(\{ 3 \}\right) = \frac{1}{4}
    = \frac{1}{2} \cdot \frac{1}{2} = \pr(Y) \cdot \pr(Z),
\end{align*}
$X, Y, Z$ are three random events such that any pair of them is independent.
However, since
\[
  \pr(X \cap Y \cap Z) = \pr(\emptyset) = 0
    < \left(\frac{1}{2}\right)^3 = \pr(X) \cdot \pr(Y) \cdot \pr(Z)
\]
$X, Y$ and $Z$ are not mutually independent.

\paragraph{Exercise 1.22}
\begin{enumerate}
  \item[(a)] If each $1 \leq i \leq n$ is chosen independently and uniformly at
  random with probability $1/2$, then each of the $2^n$ possible subsets of
  $\{ 1, ..., n \}$ is chosen with probability $2^{-n}$, giving the statement.

  \item[(b)] Let us consider $\pr\left( X \subseteq Y \right)$. $Y$ is a superset
  of $X$ if and only if all elements of $X$ are also elements of $Y$. Let $X$ be
  a particular subset of $\{ 1, ..., n \}$. In order to generate a superset of $X$
  our coin has to land heads for all elements added to $X$. Therefore the probability
  that $Y$ is a superset of the particular set $X$ is $2^{-|X|}$, and
  \[ \pr\left( X \subseteq Y \right) = \sum_{X \in \mathcal{P}\left(\{1, ..., n\}\right)} 2^{-|X|}.\]
  Since there are $\binom{n}{i}$ subsets of $\{ 1, ..., n \}$ with exactly $i$
  elements,
  \[ \pr\left( X \subseteq Y \right)
  = \sum_{X \in \mathcal{P}\left(\{1, ..., n\}\right)} 2^{-|X|}
  = \sum_{i=0}^n \binom{n}{i} \cdot \left(\frac{1}{2}\right)^{i}
  = \left( \frac{1}{2} + 1 \right)^n
  = \left( \frac{3}{2} \right)^n,\]
  where we used the binomial formula to achieve the third equation.
  Furthermore, we want to determine $\pr\left( X \cup Y = \{ 1, ..., n \}\right)$.
  Again, consider the situation there $X$ is a particular subset of $\{ 1, ..., n \}$.
  $X \cup Y = \{ 1, ..., n \}$ if and only if $Y$ is a superset of $\bar{X}$.
  The probability that $Y$ is a superset of our particular set $X$ is $2^{n-|X|}$.
  Consequently one has,
  \[ \pr\left( X \cup Y = \{ 1, ..., n \}\right)
  = \sum_{X \in \mathcal{P}\left(\{1, ..., n\}\right)} 2^{n-|X|}
  = \sum_{i=0}^n \binom{n}{i} \cdot \left(\frac{1}{2}\right)^{n-i}
  = \left( 1 + \frac{1}{2} \right)^n
  = \left( \frac{3}{2} \right)^n.\]
\end{enumerate}
